%% Shared math (response); included by supplement.tex and pre_math_report.tex.
%%%%%%%%%%%%%%%%%%%%%%%%%%%%%%%%%%%%%%%%%%%%%%%%%%%%%%%%%%%%%%%%%%%%%%
\section{Response to one seeding event}\label{sec:response}
%%%%%%%%%%%%%%%%%%%%%%%%%%%%%%%%%%%%%%%%%%%%%%%%%%%%%%%%%%%%%%%%%%%%%%

\begin{lemma}[transmission days and transmissibility; uses \ref{as:sir}]
$\Prob(D=d)=r(1-r)^d$, $d\ge0$, and
\[
\tau=\E[1-(1-p)^D]=\frac{(1-r)p}{p+r-pr}\le1-r.
\]
\end{lemma}
\begin{proof}
A node becomes infectious on some day and faces recovery that same day, so
$\Prob(D=0)=r$; each further day it transmits once and then recovers with
probability $r$, so $\Prob(D=d)=(1-r)^dr$. Then
$1-\tau=\E[(1-p)^D]=\sum_dr[(1-r)(1-p)]^d=r/(1-(1-r)(1-p))=r/(p+r-pr)$, and
$\tau=1-r/(p+r-pr)=(1-r)p/(p+r-pr)$. As $p\to1$, $\tau\to1-r$.
\end{proof}

\begin{lemma}[no secondary case; uses \ref{as:sir}]
For a seeded node with $k$ susceptible partners,
$\Prob(\text{no secondary case})=\E[(1-p)^{kD}]=r/(1-(1-r)(1-p)^k)$.
\end{lemma}
\begin{proof}
$\sum_dr(1-r)^d(1-p)^{kd}=r/(1-(1-r)(1-p)^k)$. Note it differs from $(1-\tau)^k$
because the $k$ edges share the same $D$.
\end{proof}

\begin{proposition}[linear response; uses \ref{as:net}--\ref{as:sir}]\label{prop:linear}
If $R^2=\tau^2\kappa_1\kappa_2<1$, a seeding into a class-1 node of degree $k$ causes
on average $s(k)=Ak$ secondary cases, $A=\tau(1+\tau\kappa_2)/(1-R^2)$.
\end{proposition}
\begin{proof}
Let $W_2$ ($W_1$) be the mean cluster size, including the root, started by a
class-2 (class-1) node infected along an edge. By \ref{as:tree} it has $\kappa_2$
($\kappa_1$) further partners on average, each infected with probability $\tau$; by
linearity of expectation (correlations through $D$ do not affect means)
$W_2=1+\tau\kappa_2W_1$, $W_1=1+\tau\kappa_1W_2$. Hence $W_1=(1+\tau\kappa_1)/(1-R^2)$,
$W_2=(1+\tau\kappa_2)/(1-R^2)$, finite iff $R^2<1$. The seeded node infects $\tau k$
partners on average, each starting a cluster of mean $W_2$: $s(k)=\tau kW_2$.
\end{proof}

\begin{remark}[general response]
In general we write the mean response of node $i$ as $Aw_i^\eta$. The degree
model is $w=k$, $\eta=1$. Other values of $\eta$ represent nonlinear responses,
for example when importance and spreading capacity are different but
correlated attributes.
\end{remark}

